\documentclass{article}
\usepackage{import}
\import{../../template/}{article-template}
\graphicspath{{../../images}}
\addbibresource{../../template/library.bib}

\title{Fluid Dynamics}
\author{PENG}
\date{Created: 20260627, Version: 20260627}

\begin{document}
\maketitle
\tableofcontents
\section{The Physical Properties of Fluids}
\begin{definition}[martial derivative \cite{BATCHELOR_2000_AN_INTRODUCTION_TO_FLUID_DYNAMICS}]
    Rate of change at a point moving with the fluid locally:
    \begin{equation}
        \frac{D}{Dt} = \frac{\partial}{\partial t} + \boldsymbol{u} \cdot \nabla
        = \frac{\partial}{\partial t} + [u_{x} u_{y} u_{z}]
        \begin{bmatrix}
            \frac{\partial}{\partial x} \\
            \frac{\partial}{\partial y} \\
            \frac{\partial}{\partial z}
        \end{bmatrix}
        = \frac{\partial }{\partial t} + u_{i}\frac{\partial}{\partial x_{i}}
    \end{equation}
\end{definition}

\subsection{Solids, Liquids, and Gases}
\begin{assumption}[simple fluid]
    The fluid cannot withstand any tendency by applied forces to deform it with the volume $\equiv C$.
\end{assumption}

\begin{definition}[resistance]
    The resistance cannot prevent the deformation such that resisting force vanishes with the rate of deformation.
\end{definition}

\begin{definition}[compressibility]
    gases $\gg$ liquids
\end{definition}

\begin{definition}[pressure]

\end{definition}

\subsection{The Continuum Hypothesis}

\printbibliography
\end{document}
